\section{Switching to the heteroscedastic noise model:\texorpdfstring{\\}{--}suggested sequence}
\label{app:hnm}

This appendix compares the two models for the content block of M1 and suggests when to switch. The style block
is unaffected: a style latent has no parents, so its scale cannot depend on parents and the ANM and the HNM
coincide there.

\subsection{What changes if the content block is an HNM}
\begin{itemize}[leftmargin=1.5em]
\item \textbf{Model.} The content mechanism becomes $Z_{c,i}=f_i^{(c)}(\PA)+\sigma_{c,i}^{(c)}(\PA)\,\epsilon_i$
      with $\epsilon_i$ standard normal, so the noise scale may depend on the values of the parents within a
      cell $(y,t)$, in addition to depending on $c$. In the ANM the scale depends on $c$ only.
\item \textbf{Identifiability.} Theorem~2 of \cite{ng2025} gives the same conclusion as \cref{thm:ng}: the
      graph up to isomorphism, and each latent as a function of itself and its surrounding parents. The
      assumptions are Assumption~1 and the analogue of \cref{as:2} (Assumption~3 in the paper).
\item \textbf{Sink property.} By Lemma~2, the third derivatives $\partial^3\log p/\partial Z_i^2\partial Z_j$ of a
      sink vanish only for $Z_j\notin\PA(Z_i)$. The neighbours of a sink in the Markov network are exactly its
      parents, so these entries do not vanish, and the sink no longer saves the mixed third-order entries
      of its edges. The pure third derivative $\partial^3\log p/\partial Z_i^3$ still vanishes for a sink.
\item \textbf{Length of the vector.} In Assumption~3 the mixed third-order group carries no sink restriction.
      With the ordered-pair reading of \cref{rem:ij}, every edge contributes two entries, so
      \begin{equation}
      L_{\mathrm{HNM}}=3n'+3E+D-|\Sk|=L_{\mathrm{ANM}}+\sum_{v\in\Sk}\deg_{\Mk_{Z'}}(v).
      \label{eq:LHNM}
      \end{equation}
      The paper gives no environment count for the HNM; \cref{eq:LHNM} is my count from the printed
      definition, under the ordered-pair reading.
\item \textbf{Estimation.} \cite{ng2025} state that the estimation method is nearly identical to that of
      the ANM. In the prior of \cref{sec:use}, the scale $\hat\sigma_i^{(u)}$ becomes a function of the
      parents, $\hat\sigma_i^{(u)}(\hat A_{i,:}\hat Z)$, in addition to $u$.
\end{itemize}

\subsection{Effect on the budget}
For the content block with $K_C=20$ values, condition (a) of \cref{prop:m1req} requires $L_c+1\le20$.
\Cref{tab:hnm} compares the two models for the content graphs of \cref{tab:examples}.

\begin{table}[ht]
\centering
\caption{Content-block length $L_c$ under the ANM (\cref{eq:L}) and the HNM (\cref{eq:LHNM}), and whether
$L_c+1\le K_C=20$.}
\label{tab:hnm}
\begin{tabular}{@{}lcccccc@{}}
\toprule
content graph $\Gc$ & $n_c$ & sinks (degree) & $L_{\mathrm{ANM}}$ & $L_{\mathrm{HNM}}$ & ANM fits & HNM fits\\
\midrule
empty                              & 9 & 9 (all 0) & 18 & 18 & yes & yes\\
collider $Z_1\to Z_2\leftarrow Z_3$ & 3 & 1 (2) & 16 & 18 & yes & yes (barely)\\
chain, 3 nodes                     & 3 & 1 (1) & 13 & 14 & yes & yes\\
chain, 4 nodes                     & 4 & 1 (1) & 19 & 20 & yes (barely) & no\\
chain, 5 nodes                     & 5 & 1 (1) & 25 & 26 & no & no\\
\bottomrule
\end{tabular}
\end{table}

\subsection{Suggested sequence}
\begin{enumerate}[label=\arabic*.,leftmargin=2em]
\item \textbf{Fit the ANM first.} Use the prior of \cref{sec:use}: the content scale depends on $c=(y,t)$, the
      style scale on $r=(t,e)$, and neither depends on parent values.
\item \textbf{Diagnose ANM misspecification.} Three checks, all heuristic.
      (a) Compare the held-out ELBO of the ANM with that of an HNM prior whose scale depends on the estimated
      parents; a more flexible model always fits at least as well on the training data, so use held-out data or
      a penalty.
      (b) Within each cell $(y,t)$, regress the squared residuals of each estimated content latent on the
      estimated values of its parents; a clear dependence indicates parent-dependent scale.
      (c) Check whether the recovered graph is stable across random seeds and between the two priors.
\item \textbf{Switch only if needed and affordable.} Move to the HNM if the HNM is clearly better in step 2 and
      the budget allows it, i.e.\ $L_{\mathrm{HNM}}+1\le K_C$ for the chosen content graph
      (\cref{tab:hnm}); otherwise keep the ANM, or reduce $n_c$ so that the HNM fits.
\item \textbf{Leave the style block unchanged.} It is isolated, so the two models coincide there.
\end{enumerate}

\subsection{Why the ANM first}
\begin{itemize}[leftmargin=1.5em]
\item \textbf{Smaller budget.} By \cref{eq:LHNM}, the HNM adds the total degree of the sinks to the length of
      $\vect w$. With about twenty $(y,t)$ values this is decisive for the four-node chain, which fits
      under the ANM only barely and not at all under the HNM.
\item \textbf{The ANM already allows what is most likely to matter.} The scale may change with the species and
      the time of day. What the ANM excludes is dependence of the scale on the values of the other content
      latents within one cell, a second-order effect.
\item \textbf{A cleaner sink test.} Under the ANM every third derivative $\partial^3\log p/\partial Z_i^2\partial Z_j$
      of a sink vanishes, whereas under the HNM this holds only for non-parents.
\item \textbf{A simpler prior.} The scale is a per-cell parameter and not a function of the parents, so fewer
      parameters must be estimated per environment.
\item \textbf{The need can be diagnosed afterwards.} If the residual-scale check finds no dependence on the
      parents, the switch is unnecessary; hence ``HNM later, if at all''.
\end{itemize}
